\documentclass[12pt,a4paper]{article}

\usepackage[utf8]{inputenc}
\usepackage[T1]{fontenc}
\usepackage[spanish,es-nodecimaldot]{babel}
\usepackage{amsmath,amssymb,amsthm}
\usepackage{geometry}
\usepackage{microtype}

\geometry{margin=2.5cm}

\title{Apuntes de matemáticas}
\author{}
\date{}

\begin{document}

\maketitle

\section{Introducción}

Este documento sirve como base para escribir matemáticas en \LaTeX. Las
expresiones breves pueden escribirse dentro del texto, por ejemplo,
$a^2+b^2=c^2$.

Para ecuaciones centradas y numeradas se puede usar el entorno
\texttt{equation}:

\begin{equation}
	e^{i\pi}+1=0.
\end{equation}

Las ecuaciones sin numerar se escriben con corchetes:

\[
	\int_0^1 x^2\,dx = \frac{1}{3}.
\]

Para varias líneas alineadas:

\begin{align}
	f(x) &= x^2 + 2x + 1 \\
		 &= (x+1)^2.
\end{align}

También se pueden definir teoremas y demostraciones:

\newtheorem{teorema}{Teorema}

\begin{teorema}
	Para todo $x,y\in\mathbb{R}$ se cumple $(x+y)^2=x^2+2xy+y^2$.
\end{teorema}

\begin{proof}
	Se desarrolla el producto: $(x+y)(x+y)=x^2+xy+yx+y^2$.
\end{proof}

\end{document}
